\section{Generating functions}

\begin{recognise}
The $n$th term of a sequence depends on sums of products of earlier terms, which is a convolution. A count splits an object into an ordered sequence of parts, each part counted by a known sequence. A recurrence has coefficients that are polynomials in $n$. A sum of binomial coefficients runs over every third index, or over an arithmetic progression of indices. A problem asks about the coefficients of a power of a polynomial. The seminar's eleventh set accepts only this method, and its problems are the ladder.
\end{recognise}

\begin{definition}
The ordinary generating function of a sequence $(a_n)_{n\ge0}$ is the formal power series $A(x)=\sum_{n\ge0}a_nx^n$; the exponential generating function is $\hat A(x)=\sum_{n\ge0}a_nx^n/n!$. Formal series are added and multiplied coefficientwise and by the Cauchy product; a series with nonzero constant term has a reciprocal; no convergence is involved until a number is substituted for $x$. We write $[x^n]A$ for the coefficient of $x^n$.
\end{definition}

\begin{theorem}[Products are convolutions]
For ordinary and for exponential generating functions respectively,
\[
[x^n]\,A(x)B(x)=\sum_{k=0}^na_kb_{n-k},\qquad [x^n/n!]\,\hat A(x)\hat B(x)=\sum_{k=0}^n\binom nka_kb_{n-k}.
\]
Hence: if an object of size $n$ is an ordered pair of a first part of size $k$ counted by $a_k$ and a second of size $n-k$ counted by $b_{n-k}$, the objects are counted by $AB$; if the parts are built on labelled elements and the labels are distributed between the parts, the objects are counted by $\hat A\hat B$. An ordered sequence of any number of parts, each counted by $C$ with $c_0=0$, is counted by $1/(1-C(x))$.
\end{theorem}

\begin{proof}
Multiply out: $x^k\cdot x^{n-k}=x^n$. For the exponential version, $\frac{x^k}{k!}\cdot\frac{x^{n-k}}{(n-k)!}=\binom nk\frac{x^n}{n!}$, and $\binom nk$ is the number of ways to choose which labels go to the first part. The sequence statement is $\sum_{j\ge0}C(x)^j$, a geometric series that makes sense formally because $c_0=0$ means $C^j$ starts at $x^j$ and each coefficient is a finite sum.
\end{proof}

\begin{theorem}[Linear recurrences and rational functions]
A sequence satisfies $a_{n+k}=c_{k-1}a_{n+k-1}+\dots+c_0a_n$ for all $n\ge0$ if and only if $Q(x)A(x)$ is a polynomial of degree less than $k$, where $Q(x)=1-c_{k-1}x-\dots-c_0x^k$ is the reversed characteristic polynomial. Partial fractions then give the closed form: if $Q(x)=\prod_i(1-r_ix)^{m_i}$ then $A(x)=\sum_i\sum_{j=1}^{m_i}\dfrac{\alpha_{ij}}{(1-r_ix)^j}$ and
\[
[x^n]\frac{1}{(1-rx)^j}=\binom{n+j-1}{j-1}r^n,
\]
a polynomial in $n$ of degree $j-1$ times $r^n$.
\end{theorem}

\begin{proof}
The coefficient of $x^{n+k}$ in $Q(x)A(x)$ is $a_{n+k}-c_{k-1}a_{n+k-1}-\dots-c_0a_n$, which vanishes for all $n\ge0$ exactly when the recurrence holds; so $QA$ has no terms beyond $x^{k-1}$. The coefficient formula is the binomial series $(1-y)^{-j}=\sum_n\binom{n+j-1}{j-1}y^n$, proved by differentiating the geometric series $j-1$ times, with $y=rx$.
\end{proof}

A recurrence with polynomial coefficients, like $a_{n+1}=(n+1)a_n+\dots$, becomes a differential equation for $A$, since $xA'(x)=\sum na_nx^n$; that is when the exponential generating function usually gives the simpler equation.

\begin{theorem}[Roots-of-unity filter]
Let $\omega=e^{2\pi i/k}$ and let $A(x)$ be a polynomial, or a series converging at the points $\omega^j$. Then for any residue $r$,
\[
\sum_{n\equiv r\ (\mathrm{mod}\ k)}a_n=\frac1k\sum_{j=0}^{k-1}\omega^{-jr}A(\omega^j).
\]
\end{theorem}

\begin{proof}
$\sum_{j=0}^{k-1}\omega^{j(n-r)}$ is $k$ when $k\mid n-r$ and otherwise a geometric sum $\frac{\omega^{k(n-r)}-1}{\omega^{n-r}-1}=0$. Multiply $a_n$ by this, sum over $n$, and swap the finite sum over $j$ with the sum over $n$.
\end{proof}

For $k=2$ this is $\sum_{n\text{ even}}a_n=\frac{A(1)+A(-1)}2$. For $k=3$ with $\omega^2+\omega+1=0$ the values $A(\omega)$ and $A(\omega^2)$ are conjugate when $A$ has real coefficients, so the filter is $\frac13\bigl(A(1)+2\operatorname{Re}\,\omega^{-r}A(\omega)\bigr)$.

\begin{lemma}[Standard series]
$\dfrac1{1-x}=\sum x^n$; $\dfrac{1}{(1-x)^{m}}=\sum_n\binom{n+m-1}{m-1}x^n$; $(1+x)^n=\sum_k\binom nkx^k$; $e^x=\sum x^n/n!$ is the exponential generating function of the all-ones sequence, and $e^x-1$ of the sequence that is $0$ at $n=0$ and $1$ after; $\dfrac{1-\sqrt{1-4x}}{2x}$ generates the Catalan numbers; $\dfrac{x}{1-x-x^2}$ generates the Fibonacci numbers with $F_0=0$, $F_1=1$.
\end{lemma}

\begin{proof}
The first three are the geometric and binomial series. The exponential ones are the definition. The Fibonacci one is the previous theorem with $Q=1-x-x^2$ and $QA=x$. For Catalan, the recurrence $C_{n+1}=\sum_kC_kC_{n-k}$ says $C(x)-1=xC(x)^2$, and the quadratic formula gives the root with $C(0)=1$.
\end{proof}

\begin{firstmoves}
Write the series. Multiply the recurrence by $x^{n+k}$ and sum over $n$; collect $A(x)$ and the boundary terms. Decompose the object into parts and multiply the parts' series. Find the closed form and do partial fractions. For every $k$th coefficient, evaluate at the $k$th roots of unity. For labelled objects, use exponential generating functions. When a closed form is not available, extract what is asked (a sign, a residue, a growth rate) directly from the functional equation for $A$.
\end{firstmoves}

\begin{example}
Call a string over the alphabet $\{A,B,C\}$ tidy if every maximal run of consecutive $A$'s has even length. Let $t_n$ be the number of tidy strings of length $n$. Find a recurrence for $t_n$ and a closed form.
\end{example}

\begin{proof}[Solution]
A tidy string is an alternation: a run of $A$'s of even length (possibly empty), then a letter $B$ or $C$, then an even run of $A$'s, then a letter $B$ or $C$, and so on, ending with an even run. Even runs are counted by $E(x)=1+x^2+x^4+\dots=1/(1-x^2)$, and a letter $B$ or $C$ by $2x$. A tidy string is an even run followed by an ordered sequence of blocks, each block a letter followed by an even run, so by the convolution theorem
\[
T(x)=E(x)\cdot\frac{1}{1-2xE(x)}=\frac{1}{1-x^2}\cdot\frac{1}{1-\frac{2x}{1-x^2}}=\frac{1}{1-2x-x^2}.
\]
Reading the denominator as a reversed characteristic polynomial, $t_n=2t_{n-1}+t_{n-2}$ for $n\ge2$, with $t_0=1$ and $t_1=2$. Check: $t_2=5$, and the tidy strings of length $2$ are $AA$ and the four strings with no $A$.

The characteristic roots are $1\pm\sqrt2$, so $t_n=\alpha(1+\sqrt2)^n+\beta(1-\sqrt2)^n$; from $t_0=1$ and $t_1=2$,
\[
t_n=\frac{(1+\sqrt2)^{n+1}-(1-\sqrt2)^{n+1}}{2\sqrt2}.
\]
At $n=0$ this is $2\sqrt2/(2\sqrt2)=1$ and at $n=1$ it is $4\sqrt2/(2\sqrt2)=2$. The recurrence has a direct explanation too: a tidy string ends in $B$ or $C$ preceded by a tidy string of length $n-1$, or in $AA$ preceded by a tidy string of length $n-2$; the generating function found it without that insight, and it also proves the decomposition is exhaustive and unambiguous.
\end{proof}

\begin{example}
Let $n\ge1$. Find the number of pairs $(A,B)$ of subsets of $\{1,\dots,n\}$ with $A\subseteq B$ and $|B\setminus A|$ divisible by $3$.
\end{example}

\begin{proof}[Solution]
As in the counting section, each element is outside $B$, in $B\setminus A$, or in $A$, and marking the middle state with $x$ gives the polynomial $(2+x)^n=\sum_k\binom nk2^{n-k}x^k$, whose coefficient of $x^k$ counts pairs with $|B\setminus A|=k$. The filter with $k=3$, $r=0$ and $\omega=e^{2\pi i/3}$ gives
\[
N=\frac13\Bigl(3^n+(2+\omega)^n+(2+\omega^2)^n\Bigr).
\]
Now $2+\omega=\tfrac32+\tfrac{\sqrt3}2i=\sqrt3\,e^{i\pi/6}$ and $2+\omega^2$ is its conjugate, so
\[
N=\frac13\Bigl(3^n+2\cdot3^{n/2}\cos\frac{n\pi}6\Bigr).
\]
Check at $n=1$: $\frac13(3+2\sqrt3\cdot\frac{\sqrt3}2)=2$, the pairs $(\emptyset,\emptyset)$ and $(\{1\},\{1\})$. At $n=2$: $\frac13(9+2\cdot3\cdot\frac12)=4$, the four pairs with $A=B$, since $|B\setminus A|=3$ is impossible and $|B\setminus A|\in\{1,2\}$ is excluded. At $n=3$: $\frac13(27+0)=9$, which is the $8$ pairs with $A=B$ plus $(\emptyset,\{1,2,3\})$.

The parity version, $|B\setminus A|$ odd, is the same filter with $k=2$ and gives $\frac12(3^n-1)$; the involution in the counting section proved that without complex numbers, and no such involution is in sight for divisibility by $3$, which is what the roots of unity are for.
\end{proof}

\begin{trap}
Convergence is irrelevant for a formal power series and relevant the moment a number is substituted: the roots-of-unity filter on an infinite series needs the series to converge at those points, and a closed form like $1/(1-2x)$ represents the sequence $2^n$ for $|x|<\tfrac12$ and nothing at $x=1$. A recurrence that starts late, so the first few coefficients need their own correction terms in $QA$; always compute $QA$ to the order $x^{k-1}$ explicitly rather than assuming it is $a_0$. Ordinary and exponential series mixed in one product. Partial fractions with a repeated root done as if the root were simple. A decomposition into parts that is not unique, so the product overcounts; the tidy-string decomposition was unambiguous because every block ends with the first non-$A$ letter.
\end{trap}

\begin{ladder}
\item \seminar{PS 11 \#2}. A binomial-type identity for falling factorials; prove it.
\item \seminar{PS 11 \#3}. Subsets counted by size modulo three; a closed form.
\item \seminar{PS 11 \#7}. Coefficients of a power of a short polynomial; a property of them.
\item \putnam{2017 B3}. A power series with zero-one coefficients built from a recurrence; irrationality of its value.
\item \putnam{2023 A5}. A signed sum indexed by digit counts; its complex roots.
\item \putnam{2020 B4}. A weighted average over lattice paths; evaluate.
\item \putnam{2024 A6}. A matrix built from the coefficients of a series; its determinant.
\end{ladder}

\begin{hides}
Inside recurrences and counting, as the method that does the bookkeeping (\putnam{2020 A5}, \putnam{2024 B5}, \seminar{PS 4 \#12}). Inside number theory, where coefficients are read modulo a prime or a power of a polynomial is expanded modulo $2$ (\putnam{2022 B1}, \seminar{PS 9 \#2}). Inside probability, where the generating function of a distribution gives its moments and the probability of a sum (\putnam{2020 B4}, \seminar{PS 6 \#12}). Inside sums, where a sum of binomial coefficients with a pattern is a coefficient or a value of a known series (\putnam{2020 A2}). Inside linear algebra, where Hankel determinants of a series encode a continued fraction (\putnam{2024 A6}, \seminar{PS 7 \#24}).
\end{hides}
